% Author-team research draft. Not yet integrated, compiled, or independently
% challenged. Labels refer to the complete R73 manuscript and supplement.
\section{Joint contraction and bounded positive realizations}
\label{sec:r74-joint-contraction}

The reversible classification has a complementary dissipative regime. The
condition below is imposed on all command products through one supplied norm,
and the construction gives one machine for all words. A separate realization
for each command would not establish that conclusion. Invariant-polytope
positive realization is classical \cite{Heller65,MonrasWinter14}; the point of
the statement is an explicit sufficient certificate with bounded probability
decoders, a prescribed pure start, and a stated number of live labels.

Let $A$ and $C$ be nonempty finite sets, $a=|A|$, and let
\begin{equation}\label{eq:r74-contracting-response}
 b(w,c)=p_c+\ell_c(x_w),\qquad
 x_{wa}=F_ax_w,\qquad x_\varnothing=x_0\in\R^n,
\end{equation}
where $F_a$ are linear, $p\in\Delta(C)$, $\sum_c\ell_c=0$, and the displayed
values are probabilities for every word. For a norm $\|\cdot\|$ write
$\|\ell\|_*$ for its dual norm. Suppose we are given $0<\rho<1$ and $r>0$
such that
\begin{equation}\label{eq:r74-contraction-certificate}
 \|F_ax\|\le\rho\|x\|\quad(a\in A),\qquad
 r\|\ell_c\|_*\le p_c\quad(c\in C).
\end{equation}
In particular, strictly positive $p_c$ give such an $r$ for every finite
collection of output functionals. Zero components are allowed only when the
corresponding functional vanishes. The second condition certifies that the
entire ball $rB$, not only the actual orbit, has legal decoder rows.

\begin{theorem}[A common contracting realization]\label{thm:r74-contraction}
Under \eqref{eq:r74-contraction-certificate}, the response
\eqref{eq:r74-contracting-response} has an exact stochastic realization from
a prescribed pure label, valid for every command word, with at most
\begin{equation}\label{eq:r74-contraction-count}
 \sum_{j=0}^{T-1}a^j+
 \left\lceil\left(1+\frac{2}{1-\rho}\right)^n\right\rceil
 \quad\hbox{live labels},
\end{equation}
where $T\ge1$ may be any integer satisfying
$\rho^{T-1}\|x_0\|\le r$. For $x_0\ne0$ one may take
\[
 T=1+\max\left\{0,
 \left\lceil\frac{\log(\|x_0\|/r)}{\log(1/\rho)}\right\rceil\right\};
\]
for $x_0=0$, the response is constant and one label suffices. No minimality or
computational-efficiency assertion is made for this upper bound.
\end{theorem}
\begin{proof}
Put $\varepsilon=1-\rho$. Choose a maximal $\varepsilon$-separated finite
set $V$ in the unit ball $B$, where distinct selected points have distance
greater than $\varepsilon$. The open balls of radius $\varepsilon/2$ about
these points are pairwise disjoint and lie in $(1+\varepsilon/2)B$.
Lebesgue volume, which scales by the $n$th power, gives
\[
 |V|\le (1+2/\varepsilon)^n.
\]
This bound also proves that the greedy selection stops; maximality says
that $V$ is an $\varepsilon$-net of $B$. If $P=\operatorname{conv}(V)$,
then $\rho B\subseteq P\subseteq B$. Indeed, for every real linear
functional $f$, choose a maximizer $x\in B$ and a point $v\in V$ at
distance at most $\varepsilon$ from $x$. Symmetry of the unit ball gives
\[
 \max_{v\in V} f(v)\ge f(x)-\varepsilon\|f\|_*
 =(1-\varepsilon)\|f\|_*.
\]
The separating-hyperplane characterization of convex containment proves
$\rho B\subseteq P$. Consequently
\[
 F_a(rP)\subseteq\rho rB\subseteq rP
 \quad(a\in A).
\]

Use one recurrent label for each $v\in V$, representing the vector $rv$.
For every such label and command choose nonnegative barycentric coefficients
of $F_a(rv)$ in $rV$; they sum to one and define a row of $Q_a$. Give that
label decoder $p_c+\ell_c(rv)$. The decoder is nonnegative by the second
condition in \eqref{eq:r74-contraction-certificate}, and its entries sum
to one because $p$ has mass one and $\sum_c\ell_c=0$.

In addition use one transient label for each word $w$ with $|w|<T$,
representing its actual vector $x_w$. Its decoder is the actual probability
row $b(w)$, legal by hypothesis. If $|w|<T-1$, command $a$ sends it
deterministically to the label $wa$. If $|w|=T-1$, then
\[
 \|F_ax_w\|\le\rho^T\|x_0\|\le\rho r,
\]
so $F_ax_w$ is in $rP$. Choose its barycentric row in the recurrent labels.
Start at the single empty-word label. There are exactly
$\sum_{j=0}^{T-1}a^j$ such transient labels before any duplicate states
are merged. Induction on the word shows that the probability-weighted
represented vector is always $x_w$: the transition identity holds on
each individual label, and linearity preserves it under a mixture. Pairing
with all $\ell_c$ and adding $p$ gives the entire response. This proves
the count, including the pure initial label and every first command.
\end{proof}

The assumption may also be supplied as a uniform product bound. If a known
norm $\|\cdot\|_0$ and constants $K\ge1$, $0<q<1$ satisfy
\[
 \|F_w\|_0\le Kq^{|w|}\qquad(w\in A^*),
\]
choose any $q<\rho<1$ and set
\[
 \|x\|'=\sup_{w\in A^*}\rho^{-|w|}\|F_wx\|_0.
\]
This is a norm with $\|x\|_0\le\|x\|'\le K\|x\|_0$ and
$\|F_ax\|'\le\rho\|x\|'$. Hence
Theorem~\ref{thm:r74-contraction} applies with
$\|x_0\|'\le K\|x_0\|_0$ and any $r$ satisfying
$r\|\ell_c\|_{0,*}\le p_c$. This is a consequence of a supplied uniform
certificate, not a decision procedure for joint spectral radius. A common
Euclidean contraction certificate, for example, is the finite collection
$F_a^{\mathsf T}HF_a\preceq\rho^2H$ for a supplied positive definite $H$.
The actual net and barycentric rows provide finite implementation data;
their exact algebraic or sampling encodings remain distinct resources.

\section{Average mixing does not control conditional response memory}
\label{sec:r74-mixing-obstruction}

Positive realization via a common invariant polytope, including composition
of adaptive instruments, is established machinery
\cite{MonrasWinter14,Okay26}. A hypothesis forcing every normalized branch
into an interior region is stronger than fast convergence of the channel
obtained by ignoring its outcomes. The following exact example separates
these statements even when every outcome has a probability bounded away
from zero. The infinite-memory conclusion is a response-level consequence
of the inherited relative spectral obstruction; it is not a claim that
ordinary mixing theory has failed.

Fix $\sin\theta\ne0$. For a qubit put
\[
 U_\theta=\begin{pmatrix}
 \cos(\theta/2)&-\sin(\theta/2)\\
 \sin(\theta/2)&\cos(\theta/2)
 \end{pmatrix},\qquad
 E=\begin{pmatrix}3/4&0\\0&1/4\end{pmatrix},\qquad
 \tau=|0\rangle\langle0|,
 \qquad \mathcal D(X)=\tr(X)I/2.
\]
Define three linear maps
\begin{equation}\label{eq:r74-depolarizing-instrument}
 \Psi_0(X)=\tfrac18U_\theta XU_\theta^*,\quad
 \Psi_1(X)=\tfrac14\tr(EX)\tau,\quad
 \Psi_2=\mathcal D-\Psi_0-\Psi_1.
\end{equation}

\begin{theorem}[Exact memory in a uniformly positive one-dependent family]
\label{thm:r74-depolarizing-obstruction}
The maps \eqref{eq:r74-depolarizing-instrument} form a quantum
instrument. Its unconditioned channel is the completely depolarizing map
$\mathcal D$. For every density matrix, its three outcome probabilities
belong respectively to
\[
 \{1/8\},\qquad[1/16,3/16],\qquad[11/16,13/16].
\]
Started at $I/2$, its stationary outcome process is one-dependent: two
finite blocks separated by at least one unobserved outcome are independent.
Its word-probability Hankel matrix has rank exactly three. Write
$z=e^{i\theta}$. If $z$ has finite order $q$, the least classical
hidden-state count reproducing every word probability is exactly $q$ when
a mixed initial row is free, and exactly $q+1$ when the initial label must
be pure and counted. If $z$ is not a root of unity, no finite classical
model exists, even with a freely mixed start and arbitrary real stochastic
entries. All primitive superoperators are algebraic whenever
$\cos\theta$ and $\sin\theta$ are algebraic.
\end{theorem}
\begin{proof}
Use the Choi convention
$J(\Phi)=\sum_{i,j=0}^1|i\rangle\langle j|\otimes
\Phi(|i\rangle\langle j|)$. Then
\[
 J(\mathcal D)=\tfrac12I_4,
 \qquad \|J(\Psi_0)\|_{\mathrm{op}}=\tfrac14,
 \qquad \|J(\Psi_1)\|_{\mathrm{op}}=\tfrac3{16}.
\]
The first two branch Choi matrices are positive semidefinite. Subtraction
therefore gives $J(\Psi_2)\succeq I_4/16$, proving complete positivity.
The sum is trace preserving by construction. The probability bounds follow
from $\tr\Psi_0(\sigma)=1/8$ and
$\tr\Psi_1(\sigma)=\tr(E\sigma)/4\in[1/16,3/16]$.
Conjugation by $U_\theta$ depends algebraically on $\cos\theta$ and
$\sin\theta$, which proves the last assertion.

The initial state $\sigma_*=I/2$ is fixed by $\mathcal D$, giving
stationarity. To check one-dependence directly, write $\Psi_u$ for the
chronological branch product of a block $u$. Summing over one intervening
outcome gives
\[
 \sum_{s=0}^2\tr\bigl(\Psi_v\Psi_s\Psi_u(\sigma_*)\bigr)
 =\tr\bigl(\Psi_v\mathcal D\Psi_u(\sigma_*)\bigr)
 =\tr\Psi_u(\sigma_*)\,\tr\Psi_v(\sigma_*).
\]
Further unobserved outcomes leave $\mathcal D$ unchanged. Summing over
specified block events proves the claimed independence.

After outcome $1$ the normalized state is $\tau$, and that first outcome
has probability $1/8$. Repeating outcome $0$ rotates this state with
the factor $8^{-n}$. The probability of the full word $1\,0^n\,1$ is
\begin{equation}\label{eq:r74-conditioned-oscillation}
 u_n=\frac1{64}\,8^{-n}
       \left(1+\frac12\cos(n\theta)\right).
\end{equation}
All reachable subnormalized states are real symmetric matrices, a space
of real dimension three, so the word-probability Hankel rank is at most
three. For the reverse inequality take prefixes $1\,0^i$ and suffixes
$0^j\,1$, $0\le i,j\le2$. Their Hankel submatrix is $(u_{i+j})$.
Equation~\eqref{eq:r74-conditioned-oscillation} expresses it as the
product of a three-point Vandermonde matrix, a diagonal matrix with three
nonzero coefficients, and the transpose Vandermonde matrix. It is therefore
nonsingular, proving rank three.

Suppose a finite classical model exists. Write its nonnegative outcome
matrices as $T_0,T_1,T_2$, with $\sum_sT_s$ row-stochastic, and its initial
probability row as $\mu$. Then
\[
 u_n=vT_0^nd,\qquad v=\mu T_1\ge0,\quad d=T_1\one\ge0.
\]
Here $T_0$ is row-substochastic. Delete states unreachable from the support
of $v$ and states from which no positive terminal value of $d$ is reachable.
The remaining principal submatrix $S$ still realizes $u_n$. For every
remaining $i,j$, choose a positive path from $v$ to $i$, of length $a_i$
and weight $p_i>0$, and a path from $j$ to a positive terminal value, of
length $b_j$ and weight $r_j>0$. Nonnegativity gives
\[
 p_ir_j(S^n)_{ij}\le u_{a_i+n+b_j}\le C_{ij}8^{-n}.
\]
Thus $(8S)^n$ is entrywise bounded, and every eigenvalue of $S$ has modulus
at most $1/8$.

The three exponential modes of \eqref{eq:r74-conditioned-oscillation}
are $1/8$, $z/8$, and $\overline z/8$, all distinct and with nonzero
coefficients. Cayley--Hamilton and the Vandermonde independence of these
sequences force $z/8$ to be an eigenvalue of $S$. In an upper block
triangular decomposition of $S$, it belongs to an irreducible nonnegative
diagonal block $B$. Its spectral radius is exactly $1/8$. Perron
normalization by a strictly positive eigenvector therefore makes $8B$
similar to a stochastic matrix with eigenvalue $z$. This is the elementary
normalization proved before Proposition~\ref{prop:r70-relative-spectrum}.
At a largest-modulus coordinate of a corresponding complex eigenvector,
equality in the triangle inequality forces each positive successor to have
the same modulus and value multiplied by $z$. Irreducibility propagates this
equality to every coordinate, so the relation holds on every positive edge
and every directed cycle. Following positive successors
reaches a directed cycle; its length $l\ge1$ gives $z^l=1$. This
identity proves that $z$ must be a root of unity. If its order is $q$,
every such directed cycle has length divisible by $q$, so every
realization has at least $q$ labels.

For the matching upper, suppose $z$ has order $q\ge3$. In the real Bloch
plane let $P$ be the regular $q$-gon with vertices
\[
 r_j=(\cos(2\pi j/q),\sin(2\pi j/q)),\qquad 0\le j<q,
\]
and let $\sigma_j$ be the corresponding pure qubit states. The polygon
contains the disk of radius $\cos(\pi/q)\ge1/2$. The normalized branch
$0$ permutes its vertices and branch $1$ resets to $\sigma_0=\tau$.
For an arbitrary state with Bloch vector $r$ in the unit disk, write
$h=\tr\Psi_1(\sigma)\in[1/16,3/16]$. The normalized Bloch vector after
branch $2$ is
\[
 \frac{-R_\theta r/8-h r_0}{1-1/8-h},
\]
whose norm is at most $(1/8+3/16)/(11/16)=5/11<1/2$. It therefore lies
in $P$. For each input vertex and outcome, expand the normalized output in
the same vertices and multiply the barycentric coefficients by its outcome
probability. These nonnegative rows form outcome matrices $T_s$ with
$\sum_sT_s$ row-stochastic. They reproduce each subnormalized branch
state on each label. Induction proves equality of every word probability.
The uniform mixture of the vertices represents $I/2$, giving the
$q$-label upper with a free mixed start. A stationary hidden initialization
can also be used: with $T=\sum_sT_s$, every row $\mu T^n$ represents
$\mathcal D^n(I/2)=I/2$, where $\mu$ is the uniform row. A convergent
subsequence of their Cesaro averages gives $\pi T=\pi$, and $\pi$ still
represents $I/2$, so it gives the same word probabilities on $q$ labels.
If a pure start is required, add
one label representing $I/2$ and use its branch decompositions on the first
outcome; the first outcome is generated during this counted transition.
No transition returns to that label. This gives the $q+1$ upper.

It remains to exclude a pure-start realization on exactly $q$ labels.
The preceding lower then forces the irreducible block carrying $z/8$
to use every label. Perron normalization makes $8T_0$ similar to a
stochastic matrix with primitive $q$th-root eigenvalue and exactly $q$
states. The equality argument above supplies a directed $q$-cycle using
every state. An additional positive edge would close a directed cycle of
length less than $q$, which is impossible. Thus, before normalization,
$T_0$ is a weighted directed $q$-cycle. Its successive row weights are
positive. From the prescribed starting point,
\[
 p(0^n)=8^{-n}\qquad(n\ge0).
\]
For $n=1,\ldots,q$, the left side is the product of the first $n$ row
weights encountered on this cycle. Successive ratios force every weight
to be $1/8$, so $T_0=C_q/8$ after relabelling. Next,
$p(0^n1)=8^{-n}/8$ shows that every coordinate of $d=T_1\one$ is
$1/8$. This would give
\[
 p(1\,0^n\,1)
 =e_1^{\mathsf T}T_1(C_q/8)^nd
 =8^{-n}/64,
\]
contradicting \eqref{eq:r74-conditioned-oscillation}, already at $n=0$.
The pure-start minimum is therefore $q+1$.
\end{proof}

In particular, the fully algebraic choice
\[
 U_\theta=5^{-1/2}\begin{pmatrix}2&-1\\1&2\end{pmatrix},
 \qquad e^{i\theta}=(3+4i)/5,
\]
has no finite classical realization. Its superoperator coefficients are
rational. The phase is not a root of unity, since otherwise
$z+z^{-1}=6/5$ would be a rational algebraic integer and hence an integer.
Every finite-order choice, on the other hand, gives the exact counts
$q$ and $q+1$ in a fixed physical qubit with the same uniform outcome
probability bounds. The absence of hidden-state memory cannot be inferred
from predictive rank, unconditional mixing time, or pairwise correlation
range.

Stationary one-dependent processes without a finite hidden Markov model
were established earlier, notably by the finitely dependent colorings of
Holroyd and Liggett \cite{HolroydLiggett16}. We make no first-existence
claim for that phenomenon. Here the simultaneous properties are an explicit
qubit instrument, rank-three linear prediction, uniformly positive outcome
probabilities, and an exactly depolarizing unconditioned channel, together
with a sharp finite-order memory law and a separate pure-initialization
charge.

This example satisfies a stronger average-mixing property than exponential
convergence: one unobserved step erases all prior state information. The
infinite obstruction survives in a specified observed word, whose
conditional maps retain an irrational phase. It therefore does not
contradict a finite-realization theorem requiring uniform mixing of each
normalized branch. It also explains why positive single-outcome
probabilities cannot replace that hypothesis. There is no claim that an
infinite hidden-state requirement is evidence of top-journal suitability;
the example supplies a precise boundary for the constructive theorem.

\section{An inherited sharp four-period transition}
\label{sec:r74-period-four-closure}

The four-period certificate in Proposition~\ref{prop:r73-period-four}
provides an explicit finite-probe separation. Its sufficient interval was
not the sharp exact-memory interval. The latter is already a consequence
of the inherited first-sector theorem, and is recorded here to correct
the research-status distinction rather than as a new spectral result.

\begin{corollary}[Inherited complete four-period count]
\label{cor:r74-period-four-closed}
For $b_\alpha(n)=\tfrac12[1+\alpha^n\cos(\pi n/2)]$, $n\ge0$, with a
prescribed pure initial label and randomized bounded decoder, the exact
live memory is
\[
 m(b_\alpha)=
 \begin{cases}
 2,&\alpha=0,\\
 3,&0<\alpha\le1/\sqrt3,\\
 4,&1/\sqrt3<\alpha\le1.
 \end{cases}
\]
\end{corollary}
\begin{proof}
For $\alpha>0$ the three distinct modes $1,i\alpha,-i\alpha$ give the
rank lower bound three. Apply Theorem~\ref{thm:r69-sector-phase} with
$M=3$ and $\vartheta=\pi/2\le2\pi/3$; its exact threshold is
$1/(\cos(\pi/2)+\sin(\pi/2)\tan(\pi/3))=1/\sqrt3$.
Above this value three labels are impossible, while the damped four-cycle
construction of Proposition~\ref{prop:r73-period-four} works throughout
$0\le\alpha\le1$. At $\alpha=0$ the initial probability is one and all
subsequent probabilities are $1/2$, requiring and admitting two labels.
\end{proof}

The unresolved general damping problem concerns other responses and later
transition regions not covered by an existing exact theorem. In particular,
the first transition for this four-period family is closed by the earlier
first-sector result. Its explicit rational finite-gap certificate remains a
different and valid conclusion.
